\documentclass[]{../../../../tex/classes/homework}
\usepackage{../../../../tex/packages/hebrewSupport}
\usepackage{../../../../tex/packages/mathShortcuts}
\usepackage{../../../../tex/packages/theoremsSupport}


\newcommand\udv[4]{\csb{\begin{aligned}
			u &= #1  \ & dv &= #2 \\
			du &= #3 \ & v &= #4
\end{aligned}}}

\author{שחר פרץ}
\title{חדו''א 2א $\sim$ \textit{תרגיל בית 5}}
\begin{document}
	\maketitle
	\section{}
	\begin{enumerate}[(A)]
		\item נבדוק מתאי האינטגרל $\int_1^{\infty}\frac{\sinx}{x^{\ag}}$ מתכנס. 
		\begin{itemize}
			\item אם $0 < \ag \le 1$, אז $g(x) = \sinx$ רציפה בקרן ו־$f(x) = \frac{1}{x^{\ag}}$ מונוטונית וגזירה ברציפות. נוסף על כך $\int^{x}_1 \cos(t)\dt \propto \cosx$ ולכן חסומה, ו־$f(x)$ שואפת ל־$0$ באינסוף. לכן מדיריכלה האינטגרל מתכנס. 
			
			עם זאת, האינטגרל לא מתכנס בהחלט. הראנו בתרגול שהאינטגרל של $\frac{\sof{\sinx}}{x^{\ag}}$ אינו מתכנס באמצעות השוואה לטורים בהחלט, ומהשוואה (כי $t^{\ag} \ge t$ עבור $\ag \in (0, 1]$): 
			\[ \int^{x}_1\sof{\frac{\sin t}{t^{\ag}}}\dt = \int^{x}_1\frac{\sof{\sin t}}{t^{\ag}}\dt \ge \int^{x}_1 \frac{\sof{\sin t}}{t}\dt \]
			משום שאגף ימין לא מתכנס, מהשוואה כך גם אגף ימין. 
			\item אם $\ag > 1$, ראינו בתרגול שהאינטגרל מתכנס בהחלט באמצעות השוואה, ובפרט מתכנס. 
			\item אם $\ag = 0$, אז האינטגרל $ f(x) = \int^{x}_1\frac{\sin t}{t^{\ag}}\dt = \int^{\infty}_1 \cos t\dt = \sinx$ אך $\sinx$ אינו מתכנס באינסוף. סה''כ האינטגרל $f(x)$ אינו מתכבס ובפרט אינו מתכנס בתנאי. 
			\item אם $\ag < 0$, אז $f(t) = \frac{\sin t}{t^{\ag}}\dt$ בהשוואה עם $g_\eg(t) = \frac{1}{t^{-\ag - \eg}}\dt$ נקבל: 
			\[ \frac{f(x)}{g(x)} = \frac{\frac{\sin t}{t^{\ag}}}{t^{-\ag - \eg}} = \frac{\sin t}{t^{-\eg}} \rrr{t \to \infty} \infty \]
			עוד נבחין ש־: 
			\[ \int^{x}_1\dequad \ f(t)\dt = \int^{x}_1\dequad \ t^{\ag + \eg}\dt =  \]
%			\frac{t^{\ag + \eg + 1}}{\ag + \eg + 1} {- \frac{1}{\ag + \eg  + 1}} \]
			האינטגרל לעיל מתבדר אמ''מ $\ag + \eg \le 1$. תמיד נוכל למצוא $\eg$ גדול דיו כך ש־$\ag < 1 - \eg$, וסיימנו. 
			
			\end{itemize}
		\item נתבונן באינטגרל: 
			\[ \int^{\infty}_0(-1)^{\csb{x^{2}}} \]
			קל לראות שהוא לא מתכנס בהחלט כי במקרה זה הוא שווה לאינטגרל של $1$ שאיננו מתכנס. נבדוק התכנסות בתנאי. נגדיר את הפונקציה $f(t) = (-1)^{\csb{t^{2}}}$. נוכיח באינדוקציה שמתקיים $\int^{x}_0 f(t)\dt = 0$ לכל $x$ עם כך ש־$x^{2} + 0.5 \in \Nodd$, ו־$1$ לכל $x$ כך ש־$x^{2} + 0.5 \in \Neven$. הבסיס $x = \frac{1}{\sqrt 2}$ עם ריבוע $0.5$ שסה''כ אי זוגי, ואכן האינטגרל הוא $0$. יהי $x$ עם ריבוע שלם, ונסמן ב־$x' < x$  את המספר עם השורש השלם הקודם. אכן: 
			\[ \int^{x}_0\dequad \ f(t)\dt = \underbrace{\int_{0}^{x'}\dequad \ f(t)\dt}_a + \underbrace{\int_{x'}^{x}\dequad \ f(t)\dt}_b \]
			נפרק למקרים: 
			\begin{itemize}
				\item אם $x^{2} + 0.5 \in \Neven$, אז $x'^{2} + 0. 5 \in \Nodd$ ואז $a = 0$ מה.א. ובאמצעות חישוב אפשר לראות ש–$b = 1$ וסה''כ $a + b = 1$. 
				\item אם $x^{2} + 0.5 \in \Nodd$, אז $x'^{2} + 0.5 \in \Neven$ ואז $a = 1$, ובאמצעות חישוב אפשר לראות ש־$b = -1$ וסה''כ $a + b = 0$. 
			\end{itemize}
			סה''כ לאינטגרל יש ג''ח קבוע ב־$0$ וג''ח קבוע ב־$1$ ולכן הוא איננו מתכנס. 
	\end{enumerate}
	
	\section{}
	נסמן: 
	\[ F(x) = \int^{x}_0 \frac{\cos t}{1 + t}\dt \quad \quad G(x) = \int^{x}_0 \frac{\sin t}{(1 + t)^{2}}\dt \]
	נבצע אינטגרציה בחלקים. נבחין ש־: 
	\[ F(x) = \int^{x}_0 \frac{\cos t}{1 + t}\dt = \udv{\frac{1}{1 + t}}{\cos t}{-\frac{1}{(1 + t)^{2}}}{\sin t} = \frac{\sin x}{1 + x} - \int^{x}_0 \frac{- \cos t}{(1 + t)^{2}} = \frac{\sin x}{1 + x} + G(x) \]
	משום שכאשר $x \to \infty$, בבירור $\frac{\sin x}{1 + x} \to 0$, וכאשר $x \to 0$, מהצבה פשוטה נקבל $\frac{\sin x}{1 + x} \to 0$ גם כן, ולכן הביטוי זניח, ובשתי נקודות הגבול ה''בעייתיות`` נקבל $F(x) = G(x)$ (בגבול). 
	
	בגלל ש־$\int^{\infty} \frac{1}{(1 + x)^{2}}$ מתכנס ו־$\sof{\sin x}$ חסום, מאבל נקבל ש־$G(x)$ מתכנס בהחלט. עם זאת בעבור $F(x)$ מתבדר (מהשוואה לטור הרמוני, באופן דומה לדברים שנעשו בכיתה). 
	
	\section{}
	תהי $f \co [a, \infty) \to \R$ גזירה ברציפות, ונניח שהאינטגרלים $\int^{\infty}_a(f(x))^{2}$ ו־$\int^{\infty}_a(f''(x))^{2}$ מתכנסים. נוכיח ש־$\int^{\infty}_a (f'(x))^{2}\dx$ מתכנס. 
	\begin{proof}
		ראשית נבצע אינטגרציה בחלקים: 
		\[ \int^{x}_0 f'^{2}(t)\dt = \udv{f'}{f'}{f''}{f} = \underbrace{f(x) \cdot f'(x)}_{B_x} - \underbrace{\int^{x}_0 f''(t)f(t) \dt}_{A_x} \]
		נבחין שהאינטגרל באגף ימין אכן מתכנס מא''ש הממוצעים AMGM נקבל $\frac{f''(t) + f(t)}{2} \ge \sqrt[2]{f''(t)f(t)}$ כלומר: 
		\[ A_x = \int^{x}_0 f''(t)f(t)\dt \le \frac{1}{4}\int^{x}_0(f^{2}(t) + f''^{2}(t))\dt = \frac{1}{4}\csb{\int^{x}_0 f(t)^{2}\dt + \int^{x}_0 f''^{2}(t)\dt} \]
		אגב ימין בשאיפה לאינסוף הוא קומבינציה לינארית של אינטגרלים מתכנסים ולכן כאשר $x \to \infty$ האינטגרל $A_x$ מתכנס. 
		באשר ל־$B_x$, נוכל לבצע תהליך דומה באמצעות AMGM: 
		\[ B_x = f(x)f'(x) \le \frac{1}{4}\cl{f^2(x) + f'^2(x)} \]
		בבירור $f^2(x)$ מתכנס בשאיפה לאינסוף, כי האינטגרל שלו מתכנס. משום ש־$f''^2(x)$ מתכנס (כי גם כאן האינטגרל מתכנס), נפצל למקרים. 
		\begin{itemize}
			\item אם $f''^{2}(x) \to 0$, אז $f'^2(x)$ מתכנס גם הוא, ואז יש לנו קומבינציה לינארית של שתי פונקציות מתכנסות וסיימנו. 
			\item אם $f''^{2}(x) \to \ml \neq 0$, אז $f''^{2}(x)$ הוא bounded away from zero, ואז $f'^{2}(x)$ מונוטונית. במצב זה יש לנו קומבינציה לינארית של פונקציה מתכנסת עם פונקצייה מונוטונית, וממשפט וויראשטראס הראשון שראינו בחדו''א 1א, $B_x$ מתכנס במובן הרחב. 
		\end{itemize}
		עתה נניח בשלילה ש־$B_x$ אינסופי. במצב זה, האינטגרל $\int^{\infty}_1 f(x) f'(x) \dx$ בהכרח איננו מתכנס, שכן האינטגרנד איננו מתכנס. זאת בשלילה לכך שקודם לכן הראנו שהוא מתכנס. 
		
		סה''כ כאשר $x \to \infty$ הן $B_x$ והן $A_x$ מתכנסים, דהיינו גם $B_x - A_x$ מתכנס ומהשוויון לעיל סיימנו. 
	\end{proof}
	
	\stepcounter{section}
%	\section{}
%	\begin{enumerate}[(A)]
%		\item תהי $f \co [0, \infty) \to \R$ חיובית ומונוטונית יורדת, ונניח ש־$\int^{\infty}_0 f(x)\dx$ מתכנס. נוכי חשלכל $a > 0$ הטור $\sum_{n = 1}^{\infty}f(na)$ מתכנס, וכי:
%		\[ \lim_{a \to 0^{+}}\cl{a \cdot \sumninf f(na)} = \int^{\infty}_0 \dequad f(x)\dx \]
%		\begin{proof}
%			ממשפט מההרצאה, מתקיים (עבור שינוי משתנה $\phi(x) = ax$)
%			\[ \sum_{n= 2}^{x} f(na) \le \int^{x}_2\dequad \ f(ta)\dt = \int^{xa}_2 \dequad f(t)\dt \rrr{x \to \infty} L \in \R \]
%			משום שהטור פונקציה מונוטונית עולה שחסומה ע''י פונקציה מתכנסת, נסיק שהטור מתכנס (פונקציה עולה וחסומה). עתה ניתן לדבר על הערך ההתכנסות של הטור. נגדיר $g(a) = a \sum_{n = 1}^{\infty}f(an)$. 
%			
%		\end{proof}
%	\end{enumerate}
	
	\section{}
	נחשב את הגבולות של סדרות הפונקציות הבאות: 
	\begin{enumerate}[(A)]
		\item נגדיר $f_n(x) = \frac{nx}{1 + n^{2}x^{2}}$. נקבע $x \in \R$ כלשהו, ואז:
		\[ f_n(x) = \frac{nx}{1 + n^{2}x^{2}} \rrr{n \to \infty} 0 \]
		מכאן ש־$f_n \to 0$. נבדוק האם הסדרה מתכנסת במ''ש. נבחין ש־$f_n'(x)= - \frac{n(n^{2}x^{2} - 1)}{(n^{2}x^{2} + 1)^{2}}$. בהשוואה לאפס נקבל אפסים ב־$\tl x_n := \frac{1}{n}$ (ולמעשה גם בשלילי), שם ערך הפונקציה הוא מקסימלי והוא: 
		\[ f_n(\tl x_n) = \frac{n \cdot \frac{1}{n}}{1 + n^{2}\cdot \frac{1}{n^{2}}} = \frac{1}{1 + 1} = \frac{1}{2} \]
		סה''כ: 
		\[ \sup_{x \in [0, 1]}\sof{f_n(x) - f(x)} = \sof{\frac{1}{2} - 0} = \frac{1}{2} \neq 0 \]
		לכל $n \in \N$, ובפרט בשאיפה לאינסוף. לכן פונקציות אלו אינן מתכנסות במ''ש. 
		\item עבור $f_n(x) = \frac{1}{1 + x^{n}}$ בקטע $[0, 1]$. ב־$x = 0$ נקבל את הגבול $\frac{1}{1 + 0^{n}} = \frac{1}{1} = 1$. בעבור $x \neq 0$ נקבל את הגבול: 
		\[ f_n(x) = \frac{1}{1 + x^{n}} \rrr{n \to \infty} 0 \]
		(כי חזקה אוכלת הכל). סה''כ: 
		\[ f_n \to \begin{cases}
			1 & x = 0 \\
			0 & \other
		\end{cases} \]
		נבחין ש־$f_n$ פונקציות רציפות השואפות לפונקציה שאיננה רציפה. אך גבול במ''ש משמר רציפות, ולכן הן איננה מתכנסות במ''ש. 
		\item נתבונן במשפחת הפונקציות $f_n = \sqrt{n + 1}\cdot \sin^{n}x \cosx$ בכל $\R$. 
		קל לראות שכאשר $\sin x < 0$ (מתרחש במחזורים של $2\pi$) הפונקציה היא שורש (מונוטוני בכיוונו) כפול סינוס שמחליף סימן כל $n$, ובפרט זה לא מתכנס. אם $\sin x \ge 0$, אז אם $\sinx = 1$ קיבלנו $\cosx \cdot \sqrt{n + 1}$ שכאשר $n$ שואף לאינסוף מתבדר, ועבור $\sin x \in (0, 1)$ נקבל חזקה של מספר מגודל קטן מ־$1$ כפול שורש, והראשון יותר חזקה אסימפטוטית, כלומר נשאף ל־$0$. סה''כ: 
		\[ f_n \to 0_I \]
		(הפונקציה לא מוגדרת בכל $\R$, אלא רק בעבור $I= \ccb{\cl{2 \pi k, 2\pi k + 0.5\pi} \uplus \cl{2 \pi k  + 0.5\pi, 2 \pi k + \pi}\mid k \in \N}$)
		
		במקרה זה, ההתכנסות מתבצעת במ''ש כי לכל $x \in I, \eg > 0$ המ''מ $\sqrt{n + 1} \sin ^{n}x < \eg$. 
		\item נתבונן במשפחת הפונקציות $f_n(x) = \sqrt[n]{1 + x^{n}}$. בעבור $x$ קבוע, גבול ידוע מחדו''א 1א ש־$f_n(x) \to x$ ולכן ההתכנסות היא נקודתית לזהות. אין לי התחלה של מושג אם זה במ''ש או לא. 
		\item נתבונן במשפחת הפונקציות $f_n(x) = \sumnko x^{2}(1 - x^{2})^{k -1}$. נבחין ש־: 
		\[ f_n(x) = \sumnko x^{2}(1-x^{2})^{k - 1} = x^{2}\cdot \sum_{k = 0}^{n - 1}(1-x^{2})^{k - 1} = x^{2} \cdot \frac{(1 - (1-x)^{n - 1})}{1 - (1-x)} = x \cdot\cl{1-(1-x)^{n - 1}} \]
		נקבע איזשהו $x \in \R$. ואז: 
		\[ f_n(x)= x\cl{1 - (1- x)^{n - 1}} \rrr{n \to \infty} x \cdot (1 - 0) = x \]
		כי משום ש־$1 - x \in [-1, 1]$ אז $(1 - x)^{n - 1} \to 0$. מכאן שהפונקציות לעיל מתכנסות נקודתית לפונקציה $f(x) = x$. נוכיח שההתכנסות במ''ש. משום ש־$f_n(x)$ מונוטוני (עולה), וכן $f, f_n$ רציפות, וממשפט דיני $f_n$ מתכנס במ''ש ל־$f$. 
	\end{enumerate}
	
	\section{}
	\begin{enumerate}[(A)]
		\item תהי $f \co [0, \infty) \to \R$ כך ש־$f \in R[0, c]$ לכל $c > 0$. נניח כי ישנה $x_n$ מונוטונית עולה ממש, כך ש־$x_n \to \infty \land x_1 = 1$. עוד נניח שב־$X_k:=[x_k, x_{k + 1}]$ מתקיים ש־$f|_{X_k}$ שומרת סימן. נוכיח ש־$\int^{\infty}_0 f(t)\dt$ מתכנס אמ''מ $\sumninf I_n$ מתכנס כאשר $I_n = \int^{x_{n + 1}}_{x_n}f(t)\dt$. 
		\begin{proof}
			נגדיר את הפונקציה $F(x) = \int^{x}_0 f(t)\dt$. בה''כ נגדיר $x_0 =0$. עוד נגדיר את הסכום $G(k) = \sum_{n = 0}^{k} I_n$. נוכיח שלכל $k\in \N_+$ מתקיים $G(k) = F(x_k)$. באינדוקציה, בעבור $k = 0$ מתקיים $F(x_k) = 0 = G(k)$ (שני סכומים/אינטגרלים ריקים). בצעד: 
			\[ F(x_{k + 1}) = \int^{x_{k + 1}}_{0}\dequad \dequad f(t)\dt = \int_{0}^{x_k}\dequad f(t)\dt = G(k) + I_{k + 1} = I_{k + 1} = \sum_{i = 0}^{k}I_i = \sum_{i = 0}^{k + 1}I_i = G({k + 1}) \]
			\begin{itemize}
				\item אם $\int^{\infty}_0 f(t)\dt$ מתכנס, אז הראנו שגם $G(k)$ טור מתכנס. נבחין ש־$G(k) - I_0 = \sumninf I_n$ כלומר גם $\sumninf I_n$ מתכנס, לאותו הערך, כנדרש. 
				\item אם $\sumninf I_n$ מתכנס, מנימוקים דומים $G(k)$ טור מתכנס. נראה ש־$F(x)$ מתכנס (נאמר, ל־$I \in \R$). תהא סדרה $y_n$ כלשהי כך ש־$y_0 = 0, y \to \infty$. נתבונן ב־$y_n$ כללי כלשהו, ונבחין ש־$y_n \in X_{k_n}$ עבור $k_n \in \N_+$ כלשהו. משום שב־$X_{{k_n}}$ הפונקציה $f(x)$ עם אותו הסימן, בין $x_{k_n}$ ל־$x_{{k_n} + 1}$ האינטגרל $F(x)$ רק הולך וגדל (אם הסימן חיובי) או הולך וקטן (אם הסימן שלילי), ובשני המקרים נוכל לקבוע כי: 
				\[ \min(F(x_{k_n}), F(x_{k_n + 1})) \le f(y_n) \le \max(F(x_{k_n}), F(x_{k_n + 1})) \]
				ומכאן נובע ממה שהוכחנו בתחילת ההוכחה: 
				\[ \min(G({k_n}), G({k_n + 1})) \le f(y_n) \le \max(G({k_n}), G({k_n + 1})) \]
				בגלל ש־$G(k)$ מתכנס, קל לראות ש־$\min(G(k), G(k + 1))$ ו־$\max(G(k), G(k + 1))$ מתכנסים, ובפרט כל ת''ס שלהם מתכנס, ובפרט $k_n$ מתכנסת, שתיהן ל־$I$. סה''כ מסנדוויץ' $f(y_n) \to I$. מהיינה סיימנו. 
			\end{itemize}
		\end{proof}
		\item ניעזר בסעיף הקודם כדי להראות ש־$\int^{\infty}_1\frac{\sinx}{x}$ מתכנס. \begin{proof}
			עבור $x_k = \pi k$ נבחין ש־: 
			\[ \begin{cases}
				\disty I_{2k} = \int^{2\pi k + 1}_{2 \pi k}\frac{\sinx}{x} \\
				\disty I_{2k + 1} = \int^{2\pi (k + 2)}_{2 \pi (k + 1)} \frac{\sin x}{x} \dx
			\end{cases} \]
			מקיים ש־$I_{2k}$ חיוביים בעוד $I_{2k + 1}$ שליליים, ומכאן שמירת הסימן. משום שבאינטגרל על $I_k$, כל $x$ שמאסכמים עליו מקיים $\pi k < x < \pi k + 1$ אז נקבל: 
			\[ \frac{1}{\pi (k + 1)}\underbrace{\int^{\pi k}_{\pi k}\dequad \sof{\sinx}}_{C} \le \underbrace{\int^{\pi (k + 1)}_{ \pi k}\frac{\sof {\sin x}}{\sof x}\dx}_{\sof{I_k}} \le \frac{1}{\pi k}\underbrace{\int^{\pi (k + 1)}_{\pi k}\dequad \sof{\sin x}}_{C} \]
			כאשר $C$ קבוע כלשהו (שלא תלוי ב־$k$). מכאן ש־$\sof{I_k}$ סדרה מונוטונית יורדת חלש. 
			סה''כ נקבל ש־: 
			\[ \sum_{k = 1}^{\infty}I_k = \sum_{k = 1}^{\infty}(-1)^{k}\sof{I_k} \]
			מלייבניץ הטור לעיל מתכנס. מהמשפט לפני מהסעיף הקודם, מהיות $f$ שומרת סימן בתחומים $[x_k, x_{k + 1}]$, האינטגרל גם הוא מתכנס. 
		\end{proof}
	\end{enumerate}
	
	\section{}
	יהיו $f_n \co [a, b] \to \R$ $L$־ליפשיציות. נניח כי $f_n \to f$ נקודתית. 
	\begin{enumerate}[(A)]
		\item נוכיח ש־$f$ היא $L$־ליפשיץ. \begin{proof}
			נבחין שמתקיים: 
			\[ L\sof{x - y} \ge \lim_{n \to \infty}\sof{f_n(x) - f_n(y)} = \sof{\ \lim_{\mathclap{n \to \infty}}f_n(x) - \lim_{\mathclap{n \to \infty}}f_n(y)} = \sof{f(x) - f(y)} \]
			גבול משמר אי־שוויון חלש, ולכן $\sof{f(x)- f(y)} \le L\sof{x - y}$. 
		\end{proof}
		\item נוכיח ש־$f$ היא $L$־ליפשיץ. 
		\begin{proof}
			
			
			יהיו $x, y \in [a, b]$. 
			נתבונן בחלוקה כלשהי $\Pi = \{a = x_0 < x_1 < \cdots < x_n = b\}$, כך ש־$\lg(\Pi) = \frac{\eg}{L}$. נסמן $X_i = [x_i, x_{i + 1}]$. נבחין שלכל $x, y \in X_i$ מתקיים $\sof{x - y} \le \lg(\Pi) = \frac{\eg}{L}$, כלומר: 
			\[ \sof{f(x) - f(y)} \le L \sof{x - y} \le \eg \quad \land \quad \sof{f_n(x) - f_n(y)} \le L \sof{x - y} \le \eg \]
			
			לאחר הלמה לעיל, יהי $\eg > 0$. נבחין ש־: (א''ש המשולש)
			\[ \sof{f_n(x) - f(x)} \le \underbrace{\sof{f_n(x) - f_n(x_i)}}_{C} + \underbrace{\sof{f_n(x_i) - f(x_i)}}_{B} + \underbrace{\sof{f(x_i) - f(x)}}_{A} \]
			כאשר $x_i$ האיבר הכי קרוב ל־$x$ בחלוקה, שנבחר את עדינותה להיות $\lg(\Pi) \le \frac{\eg}{3L}$. מהלמה לעיל, משום ש־$x, x_i$ במרחק $\lg(\Pi)$ אחד מהשני, נקבל $A, C \le \frac{\eg}{3}$. אשר ל־$B$, מהתכנסות נקודתית אחרי $N_i$ מסוים לכל $n$ יתקיים $B \le \frac{\eg}{3}$. משום שיש כמות סופית של $x_i$־ים נוכל לסמן $N = \max_{i \in [n]}\{N_i\}$. סה''כ לכל $n > N$ בעבור אותו $x_i$ קבוע (הקרוב ביותר ל־$x$ בחלוקה) יתקיים ש־$\sof{f_n(x) - f(x)} \le 3 \cdot \frac{\eg}{3} = \eg$ לכל $x$ בקטע. סה''כ הסופרמום שואף ל־$0$ כנדרש. 
		\end{proof}
	\end{enumerate}
	
	\section{}
	תהאנה $\{f_n\}, \{g_n\}$ סדרות חסומות כך ש־$f_n \tou f, g_n \tou g$. 
	\begin{enumerate}[(A)]
		\item נניח ש־$f, g$ חסומות. נוכיח ש־$f_n \cdot g_n \tou f \cdot g$. \begin{proof}\textbf{למה: $f_n \cdot f_n \tou f \cdot f$}
			\[ \Delta(x) := \sof{(f_n\cdot f_n)(x) - (f \cdot f)(x)} = \sof{f_n(x)^{2} - f(x)^{2}} = \sof{f_n(x) - f(x)} \cdot \sof{f_n(x) + f(x)} \]
			בגבול, כאשר הפונקציות חסומות ע''י $M$ (מהנתון): 
			\[ \lim_{{n \to \infty}} \sup_{x \in I}\Delta(x) \le \lim_{{n \to \infty}} \overbrace{\sup_{x \in I}\sof{f_n(x) - f(x)}}^{\to 0} \cdot \overbrace{\sup_{x \in I} \sof{f_n(x) + f(x)}}^{\to 2M} = 0 \cdot 2M = 0 \]
			סה''כ $f_n^{2} \tou f^{2}$ כנדרש. עתה נפנה למקרה של $f, g$ כלליים. נבחין ש־: 
			\[ f \cdot g = \frac{1}{4}\cl{\cl{f + g}^{2} - \cl{f - g}^{2}} \]
			ובאופן דומה על $f_n \cdot g_n$. 	ממשפט מההרצאה, $f_n + g_n \tou f + g$ ו־$f_n - g_n \tou f - g$. מהלמה שהוכחנו $(f_n  +g_n)^{2} \tou (f + g)^{2}$ וגם $(f_n - g_n)^{2} \tou (f - g)^{2}$. סה''כ מססקווי־לינאריות של סופרמום ולינאריות הגבול, התכנסות במ''ש לינארית ולכן גם כפל ב־$\frac{1}{4}$ לא ישנה התכנסות במ''ש. סה''כ $f_n \cdot g_n \tou f \cdot g$, כנדרש.
		\end{proof}
		\item נוכיח שאין זה עובד במידה ו־$f, g$ אינן חסומות. \begin{proof}[הפרכה]
			נגדיר את שתי סדרות הפונקציות הבאות שלחלוטין לא לקחתי מהרשימות של שירי: 
			\[ f_n(x) = \begin{cases}
				\frac{1}{x} & x \neq 0 \\
				0 & x = 0
			\end{cases} \quad \quad g_n(x) = \frac{1}{n} \]
			(כלומר $f_n$ סדרה קבועה, $g_n$ פונקציה קבועה). נבחין ש־$f_n \cdot g_n = \frac{1}{nx}$ (פרט ל־$0$ שם שוויון ל־$0$). בגלל ש־$g_n \tou 0$ (התכנסות נקודתית של פונקציות קבועות היא התכנסות במ''ש) ו־$f_n(x) \tou f_1$ (סדרה קבועה מתכנסת לערך איבריה) כלומר אם המשפט היה נכון $f_n \cdot g_n \tou 0$. אך מארכימדיאניות, לכל $R$ קיים $x \in \R$ קטן דיו כך ש־$\frac{1}{nx} > R$, כלומר מהגדרת סופרמום הפונקציות לא מתכנסות במ''ש. 
		\end{proof}
	\end{enumerate}
	
	
	\ndoc
\end{document}